\documentclass[12pt]{article}
\usepackage[brazil]{babel}
\usepackage[utf8]{inputenc}
\usepackage[T1]{fontenc}
\usepackage{setspace}
\usepackage{geometry}
\usepackage{xcolor}
\usepackage{tcolorbox}
\geometry{margin=2.5cm}

% Cores
\definecolor{azul}{RGB}{0,102,204}
\definecolor{verde}{RGB}{0,153,76}
\definecolor{roxo}{RGB}{102,0,153}
\definecolor{laranja}{RGB}{255,128,0}

\title{
\textcolor{roxo}{Roteiro de Apresentação Oral}\\
\textbf{\textcolor{azul}{Tecnologias de Agricultura de Precisão: Coleta, Processamento e Análise de Dados com Arduino}}
}

\author{}
\date{}

\begin{document}

\maketitle

\onehalfspacing

\begin{tcolorbox}[colback=blue!5,colframe=azul,title=\textbf{Abertura}]
Bom dia / boa tarde.

Meu nome é \underline{\hspace{5cm}} e eu vou apresentar um projeto de \textbf{monitoramento ambiental usando Arduino}.

Esse projeto mostra como a tecnologia pode ajudar a cuidar do meio ambiente e também facilitar o aprendizado de ciências na escola.
\end{tcolorbox}

\begin{tcolorbox}[colback=green!5,colframe=verde,title=\textbf{Introdução}]
Esse trabalho foi desenvolvido para mostrar como um \textbf{sistema simples e barato} pode monitorar variáveis ambientais importantes.

O projeto usa a plataforma \textbf{Arduino} para medir temperatura, umidade do ar e umidade do solo, informações muito importantes para a \textbf{agricultura} e para o uso consciente da água.

Ele é pensado principalmente para regiões como o semiárido, onde acompanhar as condições do ambiente ajuda na sustentabilidade e na produção agrícola.
\end{tcolorbox}

\begin{tcolorbox}[colback=orange!5,colframe=laranja,title=\textbf{Objetivo do Projeto}]
O objetivo principal do projeto é mostrar que \textbf{sensores de baixo custo}, junto com conceitos básicos de estatística, podem ser usados para analisar dados ambientais em tempo real.

Além disso, o projeto busca ser uma ferramenta educativa e fácil de reproduzir, tanto para escolas quanto para pequenas propriedades rurais.
\end{tcolorbox}

\begin{tcolorbox}[colback=purple!5,colframe=roxo,title=\textbf{Como o Projeto Funciona}]
O funcionamento do sistema acontece em três etapas simples:

\textbf{1º passo -- Coleta de dados:} Sensores conectados ao Arduino medem a temperatura, a umidade do ar e a umidade do solo em intervalos regulares.

\textbf{2º passo -- Armazenamento e análise:} Os dados são armazenados e analisados usando \textbf{Python}, com cálculos como média, mediana e desvio-padrão.

\textbf{3º passo -- Interpretação:} Os resultados são organizados em tabelas e gráficos, facilitando a visualização de padrões e variações ao longo do tempo.
\end{tcolorbox}

\begin{tcolorbox}[colback=orange!5,colframe=laranja,title=\textbf{Sensores Utilizados}]
\textbf{Sensor de Temperatura e Umidade do Ar (DHT11):}
Esse sensor é responsável por medir a temperatura do ambiente e a umidade do ar. Ele é bastante utilizado em projetos por ser barato, fácil de usar e fornecer dados suficientes para análises básicas do clima.

\textbf{Sensor de Umidade do Solo (HD-38):}
Esse sensor mede a quantidade de água presente no solo. Ele funciona verificando a condutividade elétrica do solo, ajudando a identificar quando a terra está seca ou úmida, o que é muito importante para o controle da irrigação.

Esses sensores permitem acompanhar as condições ambientais em tempo real, auxiliando na tomada de decisões e no uso consciente da água.
\end{tcolorbox}

\begin{tcolorbox}[colback=blue!5,colframe=azul,title=\textbf{Resultados}]
Os resultados mostraram variações da temperatura ao longo do dia e uma maior estabilidade na umidade do solo.

O uso da estatística ajudou a identificar padrões, tendências e possíveis valores fora do normal, mostrando que o sistema é confiável e eficiente.
\end{tcolorbox}

\begin{tcolorbox}[colback=green!5,colframe=verde,title=\textbf{Conclusão}]
Com esse projeto, foi possível perceber que a união entre \textbf{Arduino, sensores e estatística} é uma solução simples e eficiente para o monitoramento ambiental.

Além de ajudar na agricultura, o sistema é uma ótima ferramenta para o \textbf{ensino de ciências}, pois permite que os alunos aprendam na prática.

O projeto mostra que tecnologia acessível pode contribuir para a sustentabilidade e para a formação científica dos estudantes.

\textbf{Muito obrigado pela atenção!}
\end{tcolorbox}

\end{document}

